% TOP_PROBLEM: {"problemId": "problem.polynomial-pivot-rule-for-the-simplex-method", "canonicalTitle": "Polynomial pivot rule for the simplex method", "releaseRank": 131, "primaryDomain": "cryptography_coding_information_optimization", "primaryDomainLabel": "Cryptography, coding, information & optimization", "displayStatus": "open", "releaseStatus": "open"}
% !TeX program = lualatex
% Definition and sources only; this file is not an LLM solution attempt.
\documentclass[11pt]{article}
\usepackage[margin=25mm]{geometry}
\usepackage{fontspec}
\setmainfont{Latin Modern Roman}
\usepackage{amsmath,amssymb,mathtools}
\usepackage{unicode-math}
\setmathfont{Latin Modern Math}
\usepackage{xurl,hyperref}
\hypersetup{colorlinks=true,urlcolor=blue}
\setcounter{secnumdepth}{0}
\setlength{\parindent}{0pt}
\setlength{\parskip}{5pt}
\setlength{\emergencystretch}{3em}
\begin{document}
\section*{\texorpdfstring{Polynomial pivot rule for the simplex method}{Polynomial pivot rule for the simplex method}}\label{131}
\subsection{Definitions and mathematical statement}
In standard form a rational linear program has $Ax=b$, $x\ge0$, and a linear objective. A feasible basis selects an invertible basic submatrix and yields a feasible basic solution. A simplex pivot exchanges one basic and one nonbasic variable according to the rule while maintaining the algorithm's feasibility and objective conventions. For bounded nondegenerate inputs, seek a rule with
\[
 \sup_{\text{allowed LPs and feasible bases}}\mathbb E[\text{pivots to an optimum}]
 \le p(m,n)
\]
for one polynomial independent of coefficient bit lengths. For a deterministic rule omit expectation. The catalogue requests a pivot-count bound and does not separately impose an efficient per-pivot selection model.
\par\smallskip\textit{Formulation and concept references:} \href{https://mathconjectures.com/conjectures/OPT-002}{[S1]}.

\subsection{Short English statement}
Is there a simplex pivot rule that always reaches an optimum after polynomially many pivots, independent of coefficient magnitudes?
\subsection{Sources}
\begin{itemize}
\item[S1]Math Conjectures, OPT-002, catalog snapshot 31 July 2026.
\url{https://mathconjectures.com/conjectures/OPT-002}.
\textit{Formulation source.}
\end{itemize}
\end{document}
